\documentclass[10pt]{amsart}

\usepackage[margin=1in]{geometry}
\usepackage[T1]{fontenc}
\usepackage{lmodern}
\usepackage{amsmath,amssymb,amsthm,mathtools}
\usepackage{microtype}
\usepackage{url}
\usepackage{enumitem}
\usepackage{alltt}
\usepackage{booktabs}
\usepackage[dvipsnames]{xcolor}
\usepackage[colorlinks=true,linkcolor=MidnightBlue,citecolor=OliveGreen,urlcolor=MidnightBlue]{hyperref}
\usepackage[capitalize,nameinlink]{cleveref}

\setlist{leftmargin=*,itemsep=2pt,topsep=4pt}
\emergencystretch=2.5em

\theoremstyle{plain}
\newtheorem{theorem}{Theorem}[section]
\newtheorem{lemma}[theorem]{Lemma}
\newtheorem{proposition}[theorem]{Proposition}
\newtheorem{corollary}[theorem]{Corollary}
\newtheorem{importedtheorem}[theorem]{Imported Theorem}
\newtheorem{citedtheorem}[theorem]{Cited Theorem}
\theoremstyle{definition}
\newtheorem{definition}[theorem]{Definition}
\newtheorem{example}[theorem]{Example}
\theoremstyle{remark}
\newtheorem{remark}[theorem]{Remark}
\crefname{importedtheorem}{Imported Theorem}{Imported Theorems}
\crefname{citedtheorem}{Cited Theorem}{Cited Theorems}
\numberwithin{equation}{section}

% The Tate--Shafarevich group (Cyrillic Sha from the AMS wncyr fonts).
\DeclareFontFamily{U}{wncy}{}
\DeclareFontShape{U}{wncy}{m}{n}{<->wncyr10}{}
\DeclareSymbolFont{mcy}{U}{wncy}{m}{n}
\DeclareMathSymbol{\Sha}{\mathord}{mcy}{"58}

\newcommand{\A}{\textup{[A]}}
\newcommand{\Oa}{\textup{[O]}}
\newcommand{\Z}{\mathbb{Z}}
\newcommand{\N}{\mathbb{N}}
\newcommand{\Q}{\mathbb{Q}}
\newcommand{\rate}{2^{O(L/\sqrt{\log L})}}
\newcommand{\nr}{\dag}% cited from the standard literature, not re-read for this note
% Lean names: typewriter, with line breaks allowed after '.' and '_'.
\DeclareUrlCommand\lean{\urlstyle{tt}\def\UrlBreaks{\do\.\do\_}\def\UrlBigBreaks{}}
% Repository names: typewriter, with line breaks allowed after '/' and '-'.
\DeclareUrlCommand\repo{\urlstyle{tt}\def\UrlBreaks{\do\/\do\-}\def\UrlBigBreaks{}}

\begin{document}

\title[Representation numbers and Tunnell's test in time $n^{o(1)}$]{Representation numbers of quadratic forms\\ and Tunnell's congruent-number test in time $n^{o(1)}$,\\ assuming the new $\#$SAT algorithms}

\author{Vaibhav Tulsyan}

\date{November 2026 (5 November 2026)}

\subjclass[2020]{11Y16, 11E25, 11G05, 11F67, 68Q25, 03B35}
\keywords{Representation numbers, ternary quadratic forms, congruent numbers, Tunnell's theorem, central $L$-values, quadratic twists, exact counting, Exponential Time Hypothesis, Lean~4}

\begin{abstract}
For a fixed positive-definite integer matrix $A$, the representation number $r_A(n)=\#\{x\in\mathbb{Z}^k : x^{\top}Ax=n\}$ is a classical object. Four ternary representation numbers decide Tunnell's congruent-number test, and through Waldspurger's theorem such numbers give central values of twisted $L$-functions. However, for a general ternary form the best deterministic bound we know for a single $r_A(n)$ is $n^{1/2+o(1)}$, and we know of no bound of the form $n^{o(1)}$, even conditionally. We present a deterministic algorithm that outputs $r_A(n)$ in time $2^{O(L/\sqrt{\log L})}=n^{o(1)}$, where $L$ is the bit length of $n$, for every fixed positive-definite $A$ of any dimension. Our approach shifts $x$ into a box of naturals by an adjugate Cauchy--Schwarz bound, splits signs so that $Q(u-B)=n$ becomes an identity between two nonnegative straight-line programs, pads the input so that all values fit, and counts the solutions with Anthropic's counting engine through our earlier network library. We prove three theorems in Lean~4 (\texttt{repNum\_TM\_rate}, \texttt{diagRep\_TM\_rate}, \texttt{tunnell\_TM\_rate}; axioms \texttt{propext}, \texttt{Classical.choice}, \texttt{Quot.sound}): a finite-alphabet Mathlib Turing machine outputs the general count, the diagonal count, and Tunnell's test bit within $\lfloor C\cdot 2^{CL/\sqrt{\log_2(L+2)}}\rfloor$ steps. On paper, Tunnell's Theorem~3 then gives $L(E_n,1)$ exactly for squarefree $n$, and a theorem of B\"ocherer and Schulze-Pillot gives central values of imaginary quadratic twists of a fixed curve of squarefree conductor, in time $n^{o(1)}$; if OpenAI's unrefereed full-BSD result holds, the counts also give the order of the Tate--Shafarevich group of $E_n$ in analytic rank~$0$. The constants are astronomical, the machines use exponential space, and everything rests on Anthropic's Lean development of its claimed refutation of the Exponential Time Hypothesis, which is under review.
\end{abstract}

\maketitle

% =====================================================================
\section{Introduction}\label{sec:intro}
% =====================================================================

How many ways can an integer be written by a quadratic form? For a positive-definite integer matrix
$A\in\Z^{k\times k}$ we write
\[
  Q_A(x)=x^{\top}Ax,\qquad r_A(n)=\#\{x\in\Z^k:\ Q_A(x)=n\}.
\]
These numbers are old. Gauss related $r_A(n)$ for $x^2+y^2+z^2$ to class numbers, and Jacobi expressed
the four-square count through divisor sums. They are also current. Tunnell~\cite{Tun83} showed that four
ternary representation numbers decide whether a squarefree $n$ is the area of a rational right triangle,
up to the Birch and Swinnerton-Dyer conjecture (BSD). His proof goes through Waldspurger's theorem, which
relates the coefficients of weight-$3/2$ theta series to central values of twisted $L$-functions. Gross~\cite{Gro87}
and B\"ocherer and Schulze-Pillot~\cite{BSP90} made this explicit for many curves.

Computing one $r_A(n)$ is not easy. Enumerating the ternary box takes $O(n^{3/2})$ steps; Tunnell wrote
that his criterion reduces the congruent-number problem ``to an algebraic computation involving
$O(D^{3/2})$ steps''~\cite[p.~325]{Tun83}. Fixing one variable and counting representations by binary
forms gives about $n^{1/2+o(1)}$ steps per $n$~\cite[\S1.3]{HTW10}. Hart, Tornar\'ia and Watkins computed
Tunnell's counts for all $n\le10^{12}$ by fast multiplication of theta series~\cite{HTW10}, which is
quasi-linear in the range but not faster for one $n$. However, for a general ternary form we know of no
deterministic algorithm for a single $r_A(n)$ in time $n^{o(1)}$, even conditionally.

In October 2026 OpenAI~\cite{OpenAI26} and Anthropic~\cite{Anthropic26} (below \Oa{} and \A) each
claimed a deterministic algorithm that counts the satisfying assignments of a $k$-CNF in subexponential
time, refuting the Exponential Time Hypothesis; both claims come with Lean~4 developments and are under
review. Rudich and Rousseau~\cite{RR26} used the engines to factor in $N^{o(1)}$, Shunia~\cite{Shu26}
proved the rate $2^{O(L/\sqrt{\log L})}$ on paper, and our factoring note~\cite{TulFac26} proved it in Lean
on \A's engine. Our class-number note~\cite{TulCN26} built a library, ArithNet, that evaluates counts of
solutions of systems with few products, and our parity note~\cite{TulPar26} extended it. We found no
derivation of representation numbers, Tunnell's counts or these $L$-values from either engine
(\cref{sec:related}).

\subsection{The problem and the running example}
We consider the \emph{representation-number problem}: for a fixed $k$ and a fixed positive-definite
$A\in\Z^{k\times k}$, given $n\in\N$ in binary, output $r_A(n)$ in binary. We also consider Tunnell's
counts. With the names of our Lean development, let
\[
  A_n=r_{x^2+2y^2+8z^2}(n),\quad B_n=r_{x^2+2y^2+32z^2}(n),\quad
  C_n=r_{8x^2+2y^2+16z^2}(n),\quad D_n=r_{8x^2+2y^2+64z^2}(n),
\]
and let \emph{Tunnell's test} be the bit $[A_n=2B_n]$ for odd $n$ and $[C_n=2D_n]$ for even $n$.

\begin{example}[Running example]\label{ex:three}
Take $n=3$. The solutions of $x^2+2y^2+8z^2=3$ are $(\pm1,\pm1,0)$, so $A_3=4$; likewise $B_3=4$. The test
asks whether $4=8$; it is false. By Tunnell's Theorem~3 (\cref{sec:tunnell}) this means
$L(E_3,1)=(A_3-2B_3)^2\beta/(16\sqrt3)=\beta/\sqrt3\approx1.5138$ for the curve $E_3:y^2=x^3-9x$, with
$\beta\approx2.62206$ a fixed period. So $E_3$ has rank $0$, and $3$ is not congruent.
\end{example}

\subsection{The naive route and the principle}
At first one might enumerate the box of candidate vectors. However, the box has $\Theta(n^{k/2})$ points.
One might instead reduce the equation to $\#$SAT and apply the engines' headline theorems. However,
Tseitin's transformation of a multiplication circuit has $\Omega(L\log L)$ clauses for every known circuit,
and \A's clause bound then gives $2^{O(L\sqrt{\log L})}$, which is superpolynomial in $n$; our factoring
note makes this point in detail~\cite[\S1]{TulFac26}. ArithNet avoids it, but it counts \emph{natural}
solutions of \emph{nonnegative} linear combinations of products. A quadratic form has integer variables
and may have negative entries.

Instead, we rely on a different principle: \emph{a positive-definite form is coercive, so after a shift every
solution is a vector of naturals in a box of side $O(n)$, and after splitting the entries of $A$ into
positive and negative parts, $Q_A(u-B)=n$ is an equality between two sums of products with natural
coefficients.} Such an equality is a straight-line program with $O(k^2)$ products, which ArithNet counts
with $O(L)$ shared bits. The values are about $n^2$, larger than the input word allows, so we pad the
input with zeros; this changes the rate by a constant factor only.

\subsection{What we prove}
In this paper we build the straight-line description of $Q_A(u-B)=n$, prove that its count is $r_A(n)$, pad
it, and prove in Lean that one Turing machine per form outputs the answer.

\begin{theorem}[\cref{sec:statement}; machine-checked]\label{thm:main}
Let $L=|\mathtt{encodeNat}\ n|$ be the length of the binary word of $n$. There are Turing machines in
Mathlib's \lean{TM2} model, with finite alphabets on all stacks, and constants $C\in\N$, such that for every
$n\in\N$ the machine halts on \lean{encodeNat}~$n$ within $\lfloor C\cdot2^{CL/\sqrt{\log_2(L+2)}}\rfloor$
steps with the following output:
\begin{enumerate}[label=(\roman*)]
\item \lean{repNum_TM_rate}: for every $k$ and every positive-definite $A\in\Z^{k\times k}$ (one machine
and one $C$ per $A$), the word \lean{encodeNat}~$r_A(n)$;
\item \lean{diagRep_TM_rate}: for every $k$ and all $a_1,\dots,a_k\ge1$, the word
\lean{encodeNat}~$\#\{x\in\Z^k:\sum_ia_ix_i^2=n\}$;
\item \lean{tunnell_TM_rate}: one machine, the one-bit word of Tunnell's test.
\end{enumerate}
\end{theorem}

\noindent Since $L=\lfloor\log_2n\rfloor+1$ for $n\ge1$, the bound is $n^{O(1/\sqrt{\log\log n})}=n^{o(1)}$,
against the trivial $O(n^{k/2})$ and, for ternary forms, $n^{1/2+o(1)}$. To our knowledge, no deterministic
algorithm with a bound of the form $n^{o(1)}$ was known for a general ternary form, even conditionally.
\Cref{sec:new} says which special forms were already within reach.

\subsection{Consequences on paper}
The theorem computes counts, and the meaning of the counts is classical. On paper, and not formalised, we
derive four consequences in \cref{sec:tunnell,sec:twists,sec:sha,sec:new}.
\begin{enumerate}[label=(\arabic*)]
\item For squarefree $n$, Tunnell's Theorem~3~\cite{Tun83} gives $L(E_n,1)$ exactly from the counts, so
$L(E_n,1)\sqrt n/\beta\in\Q$ is computed in time $n^{o(1)}$. If the test is false, $n$ is not congruent,
unconditionally. If it is true and $n\equiv1,2,3\pmod8$, the analytic rank is at least $2$, and only then is
BSD needed to conclude that $n$ is congruent. For $n\equiv5,6,7\pmod8$ the counts vanish identically.
\item For a fixed elliptic curve $E$ of squarefree conductor $N$ with $L(E,1)\ne0$, the theorem of
B\"ocherer and Schulze-Pillot~\cite{BSP90} writes $L(E,1)L(E\otimes\chi_{-D},1)\sqrt D$ as a constant times
$m_D^2$, where $m_D$ is a fixed rational combination of $O(1)$ ternary representation numbers, for the $D$ that
satisfy fixed residue conditions. So these central values are computed exactly in time $|D|^{o(1)}$.
\item If OpenAI's catalogue entry~002~\cite{OAI002} is correct (it is unrefereed), Tunnell's conjecture
$|\Sha(E_n)|=(a(n)/\sigma_0(n))^2$ holds for every odd squarefree $n$ with $L(E_n,1)\neq0$, and its analogue
for even $n$, so $|\Sha(E_n)|$ is also
computed in time $n^{o(1)}$.
\item Some special cases were already within reach: binary forms, even rank, and forms alone in their
genus reduce to factoring and known methods (for ternary genera, with class numbers, which our
class-number note computes in the same time). The new content is odd rank, where the cusp part of the
theta series carries central $L$-values.
\end{enumerate}

\subsection{Contributions}
Our contributions are:
\begin{itemize}
\item We introduce a sign-split straight-line description of $Q_A(u-B)=n$ with $16k^2+3$ steps, one product
and one chain of coefficient mass $3$ per step, and prove that its count is $r_A(n)$ once $B$ exceeds the
adjugate box bound (\cref{sec:method}).
\item We present a generic padding transformer for counting packages, which lets a description use values
of $2L+O(1)$ bits on an input of $L$ bits at a constant-factor cost in the exponent (\cref{sec:padding}).
\item We prove the three theorems of \cref{thm:main} in Lean and list what a reader must trust
(\cref{sec:statement,sec:formal}).
\item We derive on paper the consequences for Tunnell's test, for central values of imaginary quadratic twists
and, conditionally, for Tate--Shafarevich groups, and we say which special forms were already within reach
(\cref{sec:tunnell,sec:twists,sec:sha,sec:new}).
\end{itemize}
The method is simple: a box bound, a shift, a sign split and a padding step. The network library and the
machines are those of our class-number note~\cite{TulCN26}; what we add is the description, its counting
proof, the padding, and the bridge to the number theory.

\subsection{Status of each claim}\label{sec:status}
We separate four kinds of claim.
\emph{Machine-checked in Lean}: the three theorems of \cref{thm:main}, with the axioms \lean{propext},
\lean{Classical.choice} and \lean{Quot.sound} only, and every lemma whose Lean name we give.
\emph{Proved on paper only}: everything in \cref{sec:tunnell,sec:twists,sec:sha,sec:new}, which connects the
counts to congruent numbers, $L$-values and Tate--Shafarevich groups; the reduction of forms with odd
cross coefficients to \cref{thm:main} by doubling (\cref{rem:odd}); and the comparisons with earlier bounds.
\emph{Assumed}: for \cref{thm:main}, nothing beyond the three axioms, since \A's results enter as Lean
theorems of its development, which the kernel checks with ours. For the consequences on paper, the cited
theorems of the literature; BSD where we say so; and, for \cref{sec:sha} only, OpenAI's unrefereed catalogue
entry~002~\cite{OAI002}.
\emph{Checked numerically}, by short exact computations that are not part of any proof (PARI/GP~2.17.2~\cite{PARI}
and Python): the running examples; Tunnell's $L$-value formula for all $183$ squarefree $n\le300$; the
vanishing of the counts for $n\equiv5,6,7\pmod 8$ in that range; Tunnell's formula for $|\Sha|$ against
the analytic order of $\Sha$ for the $73$ squarefree $n\le300$ with $L(E_n,1)\ne0$; Gross's formula for
the curve $11a$ for all fundamental $-D$ with $D<140$ and $(\frac{-D}{11})=-1$; and Gauss's formula for
$r_3(n)$ for squarefree $5\le n<800$.

\paragraph{How we read the literature.} We read in the original: Tunnell~\cite{Tun83}, B\"ocherer and
Schulze-Pillot~\cite{BSP90}, Pal~\cite{Pal12}, Mao, Rodriguez-Villegas and Tornar\'ia~\cite{MRVT07},
Pacetti and Tornar\'ia~\cite{PT08}, Hart, Tornar\'ia and Watkins~\cite{HTW10}, Conrad~\cite{Con08} and
OpenAI's catalogue entries~\cite{OAI002,OAI006}. We could not obtain Gross's paper~\cite{Gro87} and use the
restatement in~\cite[(3.1)]{MRVT07}. We mark with~\nr{} the results that we cite from the standard
literature without having re-read the source for this note.

\subsection*{Outline}
\Cref{sec:engine} recalls \A's counting theorem and our library. \Cref{sec:method} turns a form into a
description and \cref{sec:padding} pads it. \Cref{sec:statement} quotes the statements. \Cref{sec:tunnell,sec:twists,sec:sha,sec:new}
give the consequences on paper. \Cref{sec:formal} describes the Lean development, and
\cref{sec:related,sec:limitations,sec:conclusion} cover related work, limitations and conclusions.

% =====================================================================
\section{The engine and the network library}\label{sec:engine}
% =====================================================================

In this section we recall the one theorem of \A{} that we compose with and the interface of ArithNet, both
from our earlier notes~\cite[\S2]{TulFac26} and~\cite[\S\S2--4]{TulCN26}. We use \A's Lean development at
commit \texttt{16a5ce4} (\texttt{anthropics/formal-math}, folder \texttt{eth}, Lean v4.33.1, Mathlib
\texttt{0df444a}).

A \emph{network on $[g]$ with parameters $P,B,\tau$} has, for every node $a$, settings of at most $P$ bits and
a weight $W_a$, and for every pair of nodes an interaction through $c_{a,b}$ shared bits with
$\sum_bc_{a,b}\le P$. Its value is $Z=\sum_z\prod_aW_a(z_a)\prod_{\{a,b\}}\Psi_{a,b}$ and its total is
$U=\sum_{a\ne b}c_{a,b}$. We use only label systems with $0/1$ tables, whose value is the number of
accepted labellings.

\begin{importedtheorem}[\A, Lemma~8(i) and Theorem~A(b); see~{\cite[Imported Thm.~2.2]{TulFac26}}]\label{imp:engine}
Let $K=K_\ell\ge2$ be a level of \A's schedule, $g=3K^8$, and $e_\ell\le(3x_0+6)/\sqrt{\log_2K}$ for an
absolute $x_0$. A network on $[g]$ with parameters $P,B,\tau$ and total $U$ can be evaluated within
$\bigl(K^g+2^{e_\ell K^2L'}(K(L'+B'))^{C_0K^2}2^{2gP/K}\tau\bigr)(g(P+B))^{C_1}$ operations, where
$L'=\lceil(U+gP/K)/K^2\rceil$, $B'=\lceil2gP/K\rceil+g^2B$ and $C_0,C_1$ are absolute.
\end{importedtheorem}

\noindent So a network family of size $s$ with $U=O(s)$ and $gP=O(s)$ is evaluated in $2^{O(s/\sqrt{\log s})}$ steps
at the level $K^{40}\lesssim s$ of~\cite[\S5]{TulFac26}.

\paragraph{Descriptions.} ArithNet~\cite[\S3]{TulCN26} takes a \emph{description}: variables
$x_0,x_1,\dots$ ranging over $[0,2^{n_B})$, products $x_u\cdot w$ of a variable by a variable or a constant,
chains $\sum c_t\,\mathrm{val}_t=\sum c'_t\,\mathrm{val}'_t$ with natural coefficients of mass at most $32$ per
side, comparisons, and bit patterns. Its count at size $s$ is the number of assignments in
$[0,2^{n_B})^{\mathrm{nV}}$ that satisfy everything, where the box width $n_B>s$ is fixed by the size. The
parsimony theorem (\lean{arithNet_value}) says that the network value is this count, and the parameter
bound (\lean{arithSys_U_le}) gives $U\le2^{11}c_A(D)(s+2)$ with the shape constant
$c_A(D)=64(\mathrm{nV}+\mathrm{nP}+\mathrm{nQ}+1)(\mathrm{nC}+1)$. A polynomial-time map from input words
to well-formed descriptions with bounded $c_A$ gives a counting package (\lean{ArithSpec.pkg_exists}), and
the one- and two-query machines (\lean{GR.queryM_TM}) turn a package into a Turing machine with the rate
of \cref{thm:main}. Everything in this paragraph is machine-checked in~\cite{TulCN26} and reused unchanged.

% =====================================================================
\section{From a quadratic form to a nonnegative straight-line program}\label{sec:method}
% =====================================================================

In this section we bound the solutions, shift them into the naturals, split the signs, and write the
result as a straight-line program. All of it is in the pure Mathlib modules \lean{RepQF}, \lean{SlpDesc}
and \lean{RepDesc}. Throughout, $A\in\Z^{k\times k}$ is positive definite in the sense of Mathlib's
\lean{Matrix.PosDef}, so in particular symmetric, and $Q=Q_A$.

\subsection{The box}
The first step is a bound on each coordinate that is linear in $Q(x)$ and has an integer constant.

\begin{lemma}[\lean{posDef_sq_le}, \lean{posDef_coercive}]\label{lem:box}
For every $x\in\Z^k$ and every $j$, $x_j^2\le|\mathrm{adj}(A)_{jj}|\cdot Q(x)$. Hence $x_j^2\le c(A)\,Q(x)$
with $c(A)=\sum_j|\mathrm{adj}(A)_{jj}|$.
\end{lemma}

\begin{proof}
Let $y=\mathrm{adj}(A)e_j$, so $Ay=\det(A)e_j$. Then $x^{\top}Ay=\det(A)x_j$ and
$y^{\top}Ay=\det(A)\,\mathrm{adj}(A)_{jj}$. Cauchy--Schwarz for the positive form (\lean{cs}) gives
$\det(A)^2x_j^2\le Q(x)\det(A)\,\mathrm{adj}(A)_{jj}\le Q(x)\,|\det A|\,|\mathrm{adj}(A)_{jj}|$. Since
$\det A$ is a nonzero integer, $|\det A|\le\det(A)^2$, and the claim follows.
\end{proof}

\noindent For an integer, $|x_j|\le x_j^2$, so every solution of $Q(x)=n$ lies in the box $[-cn,cn]^k$ with
$c=c(A)$ (\lean{repNum_eq_card_box}). This is crude, since $|x_j|\le\sqrt{cn}$ also holds, but the padding of
\cref{sec:padding} absorbs the extra bits.

\subsection{The shift}
Let $B=cn+1$. The map $x\mapsto u=x+B$ is a bijection between the solutions of $Q(x)=n$ and the vectors
$u\in[0,B_x)^k$ with $Q(u-B)=n$, for every box size $B_x\ge2B$, and all such $u$ lie in $[1,2B)^k$
(\lean{card_shift_two}, \lean{shift_bounds}). The machine for diagonal forms with $a_i\ge1$, and hence
Tunnell's machine, uses $c=1$ (\lean{coercive_diagonal}).

\subsection{The sign split}
Write $A=A^+-A^-$ with $A^{\pm}_{ij}=\max(\pm A_{ij},0)$. For each entry the identity
\begin{multline*}
A^+_{ij}u_iu_j+A^-_{ij}Bu_i+A^-_{ij}Bu_j+A^+_{ij}B^2\\
=A^-_{ij}u_iu_j+A^+_{ij}Bu_i+A^+_{ij}Bu_j+A^-_{ij}B^2+A_{ij}(u_i-B)(u_j-B)
\end{multline*}
holds over $\Z$. Summing over $i,j$, the left sides give a sum $S_L(u)$ of $4k^2$ monomials $c\,x_px_q$ with
natural $c$, where $x_i=u_i$ for $i<k$ and $x_k=B$, and the right sides give $S_R(u)+Q(u-B)$ with $S_R$ of
the same shape. So
\begin{equation}\label{eq:split}
  Q(u-B)=n\iff S_L(u)=n+S_R(u)\qquad\text{(\lean{monos_shift})}.
\end{equation}
Both sides of~\eqref{eq:split} are natural numbers with nonnegative coefficients; no subtraction remains.

\begin{example}[\cref{ex:three}, continued]\label{ex:three-shift}
For $A=\mathrm{diag}(1,2,8)$ only diagonal entries are nonzero and $A^-=0$, so~\eqref{eq:split} reads
$\sum_ia_i(u_i^2+B^2)=n+\sum_i2a_iBu_i$. Tunnell's machine has $c=1$, so $B=n+1=4$, and the four solutions
become $u\in\{3,5\}\times\{3,5\}\times\{4\}$. At $u=(5,5,4)$ both sides equal $379$: the left side is
$(25+16)+2(25+16)+8(16+16)$ and the right side is $3+8(5+10+32)$. The general machine of
\cref{thm:main}(i) uses $c(A)=16+8+2=26$ and $B=79$ instead.
\end{example}

\subsection{The straight-line program}
A \emph{step} computes $z:=(x_a\text{ or }0)+x_u\cdot w+\mathrm{const}$, where $w$ is a variable or a constant. It
is encoded as one product $x_u\cdot w$ and one chain $1\cdot x_a+1\cdot(x_u\cdot w)+1\cdot\mathrm{const}=1\cdot z$,
of mass $3$. A program of steps defines the description \lean{slpDesc}: the $k$ free variables, one
auxiliary variable per step, and final comparisons. When every step reads only free variables and earlier
auxiliaries, each free tuple has exactly one satisfying run (\lean{slpRun_unique}), so the count is the
number of free tuples in the box whose run satisfies the comparisons, provided that the run's values fit in
the box (\lean{slpDesc_count}).

Our program (\lean{repProg}) first sets $x_k:=B$. Then each side of~\eqref{eq:split} is an accumulator: for
the monomials $c_mx_{p_m}x_{q_m}$ it computes $y_m:=c_m\cdot x_{p_m}$ and $\mathrm{acc}_{m+1}:=\mathrm{acc}_m+x_{q_m}\cdot y_m$,
starting from $0$ on the left and from $n$ on the right. The final comparison asks that the two accumulators be
equal. There are $1+2(1+8k^2)=16k^2+3$ steps, so $\mathrm{nV}=k+16k^2+3$, $\mathrm{nP}=\mathrm{nC}=16k^2+3$ and
$\mathrm{nQ}=1$; the shape constant is $c_A=64\,(k+2(16k^2+3)+2)(16k^2+4)$ (\lean{repDesc_cA_eq}), a constant
for fixed $k$.

\begin{proposition}[\lean{repDesc_count}]\label{prop:count}
Let $x_j^2\le c\,Q(x)$ for all $x$, $|A_{ij}|\le c_m$, $cn<B$, and $K=16(k^2+1)(c_m+1)$. If $n+KB^2<B_x$ and
$n+KB^2<2^{b\ell}$ (the limb capacity), then the count of the description in $[0,B_x)$ is the number of
solutions of $Q(x)=n$ in $[-cn,cn]^k$, which is $r_A(n)$.
\end{proposition}

\noindent The proof combines the shift bijection, \eqref{eq:split}, and the bound $n+KB^2$ on every value of
the run (\lean{sumProg_bound}). Well-formedness is unconditional (\lean{repDesc_WF}).

\begin{remark}[Forms with odd cross coefficients]\label{rem:odd}
\Cref{thm:main} is about $x^{\top}Ax$ with $A$ symmetric and integral, so its cross coefficients
$2A_{ij}$ are even. An integer-valued form $Q(x)=\sum_ia_{ii}x_i^2+\sum_{i<j}a_{ij}x_ix_j$ with some $a_{ij}$
odd has $2Q=x^{\top}Gx$ for its integral Gram matrix $G$, and $r_Q(n)=r_G(2n)$. A machine for $r_Q$ then
doubles its input and runs the machine for $G$. This is on paper only; the forms in
\cref{sec:tunnell,sec:twists} that we write out have even cross coefficients.
\end{remark}

% =====================================================================
\section{Padding the input so that the values fit}\label{sec:padding}
% =====================================================================

In this section we explain why the description cannot be counted on the input word itself, and how a
generic padding step fixes this. The modules are \lean{CntPad} and \lean{RepPkg}.

The package of a description family sets the size from the input word $w$, and the box width $n_B$ grows
linearly with $|w|$. For the input \lean{encodeNat}~$n$ we have $n<2^{L}$, but the run of
\cref{prop:count} reaches $n+K(cn+1)^2$, which needs about $2L$ bits. So the box condition of
\cref{prop:count} fails on the input word. We pad instead. For a constant $c_3$, let
\[
  \mathrm{padW}_{c_3}(w)=w\,0^{|w|+c_3},
\]
so that $|\mathrm{padW}_{c_3}(w)|=2|w|+c_3$ and the little-endian value is unchanged (\lean{gofWord_padW}).
The padded package (\lean{CntPackage.padded}) answers the query of the original package on the padded word,
and keeps every bound with the constant $c_0\cdot(c_3+2)2^{\lfloor\log_2(c_3+2)\rfloor+2}$; the depth term uses
$d(2L+c_3)\le2^{\lfloor\log_2(c_3+2)\rfloor+2}\,d(L)$ (\lean{dg_pad}). The rate theorem for the padded package
(\lean{GR.queryM_TM_padded}) has its hypotheses stated on the original one.

\begin{lemma}[\lean{size_bound}, \lean{repFam_Q_pad}]\label{lem:pad}
Let $c_3=\lfloor\log_2(1+K(c+1)^2)\rfloor+1$. If $n<2^L$, then $n+K(cn+1)^2<2^{2L+c_3}$. Hence, on the padded
word, the count of the description is $\#\{x:Q(x)=n\}$ for every $n$.
\end{lemma}

\noindent The network now has size $s=2L+c_3=O(L)$, and $c_A$ is a constant, so $U,gP=O(L)$ and
\cref{imp:engine} gives $2^{O(L/\sqrt{\log L})}$. The machine for (i) asks one query (\lean{specRep}). The
machine for (iii) asks two, one for each form of the parity of $n$, and outputs $[Q_1=2Q_2]$
(\lean{specTun}); it answers $n=0$ directly, where \lean{tunnellTest}~$0$ is false. For $k=0$ the answer
$[n=0]$ needs no count (\lean{repNum_TM_zero}). The final theorems are \lean{repNum_TM_of} and
\lean{tunnell_TM_of}, instantiated in \lean{FinalRep}.

% =====================================================================
\section{What a reader must trust}\label{sec:statement}
% =====================================================================

A reader who wants to trust \cref{thm:main} must check the following definitions and nothing else of ours.
They are in the pure module \lean{RepQF}; we render Unicode symbols as mathematics. Here
\texttt{\(\cdot\sb{v}\)} is Mathlib's dot product (\lean{dotProduct}) and \texttt{*\(\sb{v}\)} its matrix--vector
product (\lean{Matrix.mulVec}).
\begin{alltt}\footnotesize
def qf \{k : \(\mathbb{N}\)\} (A : Matrix (Fin k) (Fin k) \(\mathbb{Z}\)) (x : Fin k \(\to\) \(\mathbb{Z}\)) : \(\mathbb{Z}\) := x \(\cdot\sb{v}\) A *\(\sb{v}\) x
noncomputable def repNum \{k : \(\mathbb{N}\)\} (A : Matrix (Fin k) (Fin k) \(\mathbb{Z}\)) (n : \(\mathbb{N}\)) : \(\mathbb{N}\) :=
  Set.ncard \{x : Fin k \(\to\) \(\mathbb{Z}\) | x \(\cdot\sb{v}\) A *\(\sb{v}\) x = n\}
def tForm (a b c : \(\mathbb{N}\)) : Matrix (Fin 3) (Fin 3) \(\mathbb{Z}\) := Matrix.diagonal ![(a : \(\mathbb{Z}\)), b, c]
noncomputable def tunnellTest (n : \(\mathbb{N}\)) : Bool :=
  if n % 2 = 1 then decide (repNum (tForm 1 2 8) n = 2 * repNum (tForm 1 2 32) n)
  else decide (repNum (tForm 8 2 16) n = 2 * repNum (tForm 8 2 64) n)
\end{alltt}
\lean{Set.ncard} is $0$ on infinite sets, but by \cref{lem:box} the set is finite for positive-definite $A$, and
Lean proves that it equals the finite box count (\lean{repNum_eq_card_box}). \lean{repNum} does not use
\lean{qf}; we list \lean{qf} because the coercivity bound behind \cref{lem:box} is stated with it. The final statements,
verbatim from \lean{FinalRep} up to line breaks (proofs omitted), are:
\begin{alltt}\footnotesize
theorem repNum_TM_rate (k : \(\mathbb{N}\)) (A : Matrix (Fin k) (Fin k) \(\mathbb{Z}\)) (hA : A.PosDef) :
    \(\exists\) (M : Turing.TM2ComputableAux Bool Bool) (C : \(\mathbb{N}\)), (\(\forall\) i, Finite (M.tm.\(\Gamma\) i)) \(\wedge\)
      \(\forall\) n : \(\mathbb{N}\), GraphProblems.Outputs M (Computability.encodeNat n)
        (Computability.encodeNat (repNum A n))
        \(\lfloor\)(C : \(\mathbb{R}\)) * 2 ^ ((C : \(\mathbb{R}\)) * (Computability.encodeNat n).length /
          Real.sqrt (Real.logb 2 ((Computability.encodeNat n).length + 2)))\(\rfloor\sb{+}\)
theorem diagRep_TM_rate (k : \(\mathbb{N}\)) (a : Fin k \(\to\) \(\mathbb{N}\)) (ha : \(\forall\) i, 0 < a i) :
    \(\exists\) (M : Turing.TM2ComputableAux Bool Bool) (C : \(\mathbb{N}\)), (\(\forall\) i, Finite (M.tm.\(\Gamma\) i)) \(\wedge\)
      \(\forall\) n : \(\mathbb{N}\), GraphProblems.Outputs M (Computability.encodeNat n)
        (Computability.encodeNat (Set.ncard \{x : Fin k \(\to\) \(\mathbb{Z}\) | \(\sum\) i, (a i : \(\mathbb{Z}\)) * x i ^ 2 = n\}))
        \(\lfloor\)(C : \(\mathbb{R}\)) * 2 ^ ((C : \(\mathbb{R}\)) * (Computability.encodeNat n).length /
          Real.sqrt (Real.logb 2 ((Computability.encodeNat n).length + 2)))\(\rfloor\sb{+}\)
theorem tunnell_TM_rate :
    \(\exists\) (M : Turing.TM2ComputableAux Bool Bool) (C : \(\mathbb{N}\)), (\(\forall\) i, Finite (M.tm.\(\Gamma\) i)) \(\wedge\)
      \(\forall\) n : \(\mathbb{N}\), GraphProblems.Outputs M (Computability.encodeNat n) [tunnellTest n]
        \(\lfloor\)(C : \(\mathbb{R}\)) * 2 ^ ((C : \(\mathbb{R}\)) * (Computability.encodeNat n).length /
          Real.sqrt (Real.logb 2 ((Computability.encodeNat n).length + 2)))\(\rfloor\sb{+}\)
\end{alltt}
The binary encoding is Mathlib's (\lean{Mathlib.Computability.Encoding}), little-endian with the empty word for $0$:
\begin{alltt}\footnotesize
def encodePosNum : PosNum \(\to\) List Bool
  | PosNum.one => [true]
  | PosNum.bit0 n => false :: encodePosNum n
  | PosNum.bit1 n => true :: encodePosNum n
def encodeNum : Num \(\to\) List Bool
  | Num.zero => []
  | Num.pos n => encodePosNum n
def encodeNat (n : \(\mathbb{N}\)) : List Bool := encodeNum n
\end{alltt}
\texttt{Outputs} is the one-line definition of \A's statement file \texttt{GraphProblemsStatement}, which we
also trust:
\begin{alltt}\footnotesize
def Outputs (M : TM2ComputableAux Bool Bool) (x y : List Bool) (T : \(\mathbb{N}\)) : Prop :=
  Nonempty (TM2OutputsInTime M.tm (x.map M.inputAlphabet.invFun)
    (some (y.map M.outputAlphabet.invFun)) T)
\end{alltt}
Behind these are Mathlib's \lean{Matrix.PosDef}, \lean{Matrix.diagonal}, \lean{dotProduct},
\lean{Matrix.mulVec}, \lean{Set.ncard}, \lean{TM2OutputsInTime}, \lean{Real.logb} and \lean{Real.sqrt},
from Mathlib \texttt{0df444a}. Note that \lean{tunnellTest} is defined for every $n$, while Tunnell's theorem is
about squarefree $n$; congruent numbers and $L$-functions do not occur in the statements. The constant $C$
depends on $A$ and is not tracked.

% =====================================================================
\section{Tunnell's test and exact central values}\label{sec:tunnell}
% =====================================================================

In this section we state what Tunnell's counts mean, and what follows from \cref{thm:main}. Everything here
is on paper. Let $E_n:y^2=x^3-n^2x$. A squarefree $n$ is \emph{congruent}, the area of a right triangle with
rational sides, exactly when $E_n(\Q)$ is infinite~\cite[p.~324]{Tun83}.

\paragraph{Conventions.} Tunnell writes $a(n)=\#\{2x^2+y^2+32z^2=n\}-\frac12\#\{2x^2+y^2+8z^2=n\}$ for odd $n$
and defines $b$ through the forms of the even case~\cite[pp.~325, 328]{Tun83}. With our names,
\[
  a(n)=B_n-\tfrac12A_n\quad(n\text{ odd}),\qquad b(n/2)=D_n-\tfrac12C_n\quad(n\text{ even}).
\]
Many statements of the criterion swap our $A_n$ and $B_n$, and our $C_n$ and $D_n$, and write $2A_n=B_n$; the
test is the same.

\begin{citedtheorem}[Tunnell~{\cite[Theorem~3]{Tun83}}]\label{cit:tunnell}
For squarefree $n$,
\[
  L(E_n,1)=\frac{(A_n-2B_n)^2\,\beta}{16\sqrt n}\ \ (n\text{ odd}),\qquad
  L(E_n,1)=\frac{(C_n-2D_n)^2\,\beta}{8\sqrt n}\ \ (n\text{ even}),
\]
where $\beta=\int_1^\infty dx/\sqrt{x^3-x}=2.62205\ldots$ is the real period of $E_1$ in Tunnell's
normalisation.
\end{citedtheorem}

\noindent Tunnell states it as $L(E_d,1)=a(d)^2\beta d^{-1/2}/4$ and $L(E_{2d},1)=b(d)^2\beta(2d)^{-1/2}/2$
for odd squarefree $d$; the display is the same formula in our names. We checked it against PARI/GP for all
$183$ squarefree $n\le300$ (maximal absolute difference $4.3\cdot10^{-12}$).

\begin{corollary}[on paper]\label{cor:lvalue}
For squarefree $n$, the rational number $L(E_n,1)\sqrt n/\beta$ is computed in time $n^{o(1)}$, by two runs of
the machines of \cref{thm:main}(ii).
\end{corollary}

\noindent If $n$ is not known to be squarefree, its squarefree part is found by factoring $n$, in time
$\rate$ by~\cite{Shu26,TulFac26}; $n$ is congruent exactly when its squarefree part is.

\paragraph{The congruent-number test.} The root number of $E_n$ is $+1$ for $n\equiv1,2,3\pmod8$ and $-1$
for $n\equiv5,6,7\pmod8$ (squarefree $n$). This is consistent with Tunnell's remark that $a(n)=0$ unless
$n\equiv1,3\pmod 8$ and $b(m)=0$ unless $m\equiv1,5\pmod8$~\cite[p.~329]{Tun83}.

\begin{corollary}[on paper]\label{cor:congruent}
Let $n$ be squarefree.
\begin{enumerate}[label=(\alph*)]
\item If Tunnell's test is false, then $n$ is not congruent. This is unconditional: $L(E_n,1)\neq0$ by
\cref{cit:tunnell}, so $E_n(\Q)$ is finite by Coates and Wiles~\cite{CW77}\nr, whom Tunnell quotes for this
step~\cite[p.~329]{Tun83}.
\item If $n\equiv5,6,7\pmod8$, the counts vanish identically: $A_n=B_n=0$ for odd $n$ and $C_n=D_n=0$ for even
$n$. So the test is true and carries no information.
\item If $n\equiv1,2,3\pmod8$ and the test is true, then $L(E_n,1)=0$ with root number $+1$, so the analytic
rank of $E_n$ is even and at least $2$. BSD then predicts rank at least $2$, and $n$ congruent. This is the only
case in which the test's verdict ``congruent'' needs BSD.
\end{enumerate}
\end{corollary}

\begin{proof}[Proof of (b)]
For odd $n$ in $x^2+2y^2+8z^2=n$ or $x^2+2y^2+32z^2=n$, $x$ is odd, so $x^2\equiv1$ and the left side is $1$
or $3$ modulo $8$. For even $n$ in $8x^2+2y^2+16z^2=n$ or $8x^2+2y^2+64z^2=n$, $n/2=4x^2+y^2+8z^2$ or
$4x^2+y^2+32z^2$ is odd because $n$ is squarefree, so $y$ is odd and $n/2\equiv1\pmod4$, that is $n\equiv2\pmod8$. In all other
classes the counts are $0$.
\end{proof}

\noindent For $n\equiv5,6,7\pmod8$, analytic rank $1$ already implies that $n$ is congruent, unconditionally,
by Gross--Zagier~\cite{GZ86}\nr{} and Kolyvagin~\cite{Kol88,Kol90}\nr. But our machines do not compute
$L'(E_n,1)$, so they decide nothing in these classes.

\begin{example}[\cref{ex:three}, continued; two true tests]\label{ex:true}
For $n=41\equiv1\pmod8$ we get $A_{41}=32$ and $B_{41}=16$, so the test is true and $L(E_{41},1)=0$. Here
\cref{cor:congruent}(c) needs BSD; $41$ is in fact congruent, with sides $\frac{40}{3},\frac{123}{20},\frac{881}{60}$.
For $n=34\equiv2\pmod8$ we get $C_{34}=8$ and $D_{34}=4$, and $34$ is congruent with sides
$\frac{15}{2},\frac{136}{15},\frac{353}{30}$. We found both triangles by a short search for points on $E_n$
and checked them exactly.
\end{example}

% =====================================================================
\section{Central values of imaginary quadratic twists}\label{sec:twists}
% =====================================================================

In this section we fix an elliptic curve and vary the twist. Everything here is on paper. Gross~\cite{Gro87}
proved, for prime conductor, that the central values of the imaginary quadratic twists are squares of
coefficients of a fixed combination of ternary theta series; B\"ocherer and Schulze-Pillot~\cite{BSP90}
extended this to squarefree level.

\paragraph{Setting~\cite[\S1.4, \S3]{BSP90}.} Let $N>1$ be squarefree and let $f\in S_2(\Gamma_0(N))$ be a
normalised newform, an eigenform of all Hecke operators, with sign $+$ in the functional equation of
$L(f,s)$. Let $w_q$ be its Atkin--Lehner eigenvalues, $N_1=\prod_{w_q=-1}q$ and $N_2=N/N_1$. Let $\mathcal D$ be
the definite quaternion algebra ramified exactly at the primes dividing $N_1$ and at $\infty$, let $R$ be an
Eichler order of level $N$ in $\mathcal D$ with $h$ classes represented by $y_1,\dots,y_h$, let $R_i$ be the
corresponding orders and $e_i=|R_i^\times|$, and let $L_i=\{x\in\Z+2R_i:\mathrm{tr}(x)=0\}$ with the norm form, a
positive-definite ternary lattice. Let $\varphi$ be the function on the classes attached to $f$, and
\[
  W(\varphi)=\sum_{i=1}^h\frac{\varphi(y_i)}{e_i}\,\theta(L_i,z),\qquad
  a(W;d)=\sum_{i=1}^h\frac{\varphi(y_i)}{e_i}\,r_{L_i}(d).
\]

\begin{citedtheorem}[{B\"ocherer and Schulze-Pillot~\cite[\S3, (3.5)]{BSP90}}]\label{cit:bsp}
In this setting, let $-d_0<0$ be a fundamental discriminant and $N_0=N/(N,d_0)$. Then
\[
  c_2\,d_0^{1/2}\prod_{q\mid N_0}\bigl(1+(\tfrac{-d_0}{q})w_q\bigr)\,L(f,1)\,L\bigl(f,(\tfrac{-d_0}{\cdot}),1\bigr)
  =c_1\,D_f^{(N)}(1)\,a(W;d_0N_0^2)^2,
\]
where $c_1,c_2$ are explicit constants depending only on $N$ and $D_f^{(N)}(1)$ is essentially the Petersson
norm of $f$. If $N$ is odd, $a(W;d_0N_0^2)^2=a(W;d_0)^2$.
\end{citedtheorem}

\noindent For $N$ prime this is Gross's Proposition~13.5, which we know through the restatement
of~\cite[(3.1)]{MRVT07}.

\begin{corollary}[on paper]\label{cor:twist}
Let $E/\Q$ have squarefree conductor $N$ and $L(E,1)\ne0$, and let $f$ be its newform. Let $-D<0$ be a
fundamental discriminant with $(D,N)=1$ and $(\frac{-D}{q})=w_q$ for every $q\mid N$. Then
\[
  L(E\otimes\chi_{-D},1)=\frac{\kappa_E\,m_D^2}{\sqrt D},\qquad
  \kappa_E=\frac{c_1D_f^{(N)}(1)}{2^{\omega(N)}c_2\,L(E,1)},\qquad
  m_D=\sum_{i=1}^h\frac{\varphi(y_i)}{e_i}\,r_{L_i}(DN^2),
\]
with $\kappa_E$ independent of $D$. So $m_D$, and hence $L(E\otimes\chi_{-D},1)$ up to the fixed constant
$\kappa_E$, is computed exactly in time $|D|^{o(1)}$ by $h=O(1)$ runs of the machine of \cref{thm:main}(i).
\end{corollary}

\begin{proof}
$L(E,1)\neq0$ forces the sign $+$. Since $(D,N)=1$, $N_0=N$, and each factor of the product is $2$. Divide
\cref{cit:bsp} by $2^{\omega(N)}c_2\sqrt D\,L(E,1)$. Each $r_{L_i}$ is the representation number of a fixed
positive-definite ternary form, at the argument $DN^2$ of bit length $O(L)$; if the norm form has odd cross
coefficients we use \cref{rem:odd}.
\end{proof}

\paragraph{The rational value.} With Pal's period relation~\cite[Main Result~1.1, Cor.~2.6]{Pal12}, for
squarefree $d<0$ coprime to the discriminant of $E$, $\Omega(E^d)=\tilde u\,c_\infty(E^d)\,\Omega^-(E)/\sqrt{|d|}$,
where $c_\infty$ counts real components, $\Omega^-$ is the imaginary period and $\tilde u$ is a power of $2$,
equal to $1$ when $d\equiv1\pmod4$. So $L(E^{-D},1)/\Omega(E^{-D})$ is $m_D^2$ times a constant that depends
only on $E$ and $\tilde u$. Since this ratio is rational for every twist\nr, the constant is rational, and
there are finitely many of them; each can be fixed once from one $D$ with $m_D\neq0$. Pal notes that the
formula without $\tilde u$ is misquoted in parts of the literature~\cite{Pal12}.

\paragraph{Scope.} \Cref{cor:twist} covers only the $D$ that satisfy the residue conditions. Pacetti and
Tornar\'ia describe the method of~\cite{BSP90} as covering ``only a fraction of imaginary quadratic twists
(determined by quadratic residue conditions)'', and give the composite-level version used in
computations~\cite{PT08}. For prime $N$, the excluded $D$ coprime to $N$ have root number $-1$ by the
standard formula for the root number of a twist\nr, so their central values vanish. Real twists and base curves with
$L(E,1)=0$ need the generalised theta series of~\cite{MRVT07}, with weights on residue classes; these are
counts of representations by fixed inhomogeneous quadratic polynomials, which \cref{thm:main} does not cover as
stated.

\begin{example}[The curve $11a$]\label{ex:11a}
For $E=11a$ ($y^2+y=x^3-x^2-10x-20$), $N=11$, $w_{11}=-1$ and $h=2$. Mao, Rodriguez-Villegas and Tornar\'ia write
the theta series of the eigenvector as the difference of the theta series of
\[
  Q_1=4x_1^2+11x_2^2+12x_3^2+4x_1x_3,\qquad
  Q_2=16x_1^2+15x_2^2+15x_3^2+14x_2x_3+28x_1x_3+16x_1x_2
\]
\cite[(4.1), (4.2)]{MRVT07}. Both have even cross coefficients and determinant $484$, so
\cref{thm:main}(i) applies to them directly. With $m_D=\frac12(r_{Q_1}(D)-r_{Q_2}(D))$,
\cref{tab:11a} shows that $L(E,1)\,L(E\otimes\chi_{-D},1)\sqrt D/m_D^2$ is the same number,
$0.7406174494\ldots$, for every fundamental $-D$ with $D<140$, $(\frac{-D}{11})=-1$ and $m_D\ne0$; for
$D=47,103,119$ both $m_D$ and the $L$-value vanish.
\end{example}

\begin{table}[ht]
\centering\small
\begin{tabular}{@{}rrrrr@{\qquad}rrrrr@{}}
\toprule
$D$ & $r_{Q_1}$ & $r_{Q_2}$ & $m_D$ & $L(E\otimes\chi_{-D},1)$ & $D$ & $r_{Q_1}$ & $r_{Q_2}$ & $m_D$ & $L(E\otimes\chi_{-D},1)$\\
\midrule
3 & 0 & 2 & $-1$ & 1.68450 & 59 & 8 & 6 & 1 & 0.37984\\
4 & 2 & 0 & 1 & 1.45882 & 67 & 0 & 6 & $-3$ & 3.20801\\
15 & 4 & 6 & $-1$ & 0.75333 & 71 & 16 & 18 & $-1$ & 0.34626\\
20 & 4 & 6 & $-1$ & 0.65240 & 91 & 8 & 0 & 4 & 4.89362\\
23 & 8 & 6 & 1 & 0.60837 & 103 & 12 & 12 & 0 & 0\\
31 & 8 & 6 & 1 & 0.52402 & 104 & 16 & 12 & 2 & 1.14439\\
47 & 12 & 12 & 0 & 0 & 111 & 20 & 18 & 1 & 0.27693\\
56 & 8 & 12 & $-2$ & 1.55954 & 115 & 4 & 6 & $-1$ & 0.27207\\
\bottomrule
\end{tabular}
\caption{The curve $11a$: representation numbers of the two forms of \cref{ex:11a} and the central values
of the twists, from PARI/GP. The table stops at $D=115$; the check ran to $D<140$. $L(E,1)=0.2538418609$.}\label{tab:11a}
\end{table}

% =====================================================================
\section{Tate--Shafarevich groups, conditionally}\label{sec:sha}
% =====================================================================

In this section we assume one unrefereed result and say what follows. Everything here is on paper.

OpenAI's catalogue entry~002, ``The full BSD formula from low Selmer corank''~\cite{OAI002}, consists of
preprints dated 6--7~October 2026. Its manuscript abstract states: ``We prove the full
Birch--Swinnerton-Dyer leading-term formula for every elliptic curve over $\Q$ whose full $q$-power Selmer
group has corank zero or one at some prime $q$. The analytic and Mordell--Weil ranks equal that corank, and
the Tate--Shafarevich group is finite.'' We have not checked these preprints, and they are not refereed.

If $L(E,1)\ne0$, then $E(\Q)$ and $\Sha(E)$ are finite by Kolyvagin~\cite{Kol88,Kol90}\nr, so the $q$-power
Selmer group has corank $0$ for every $q$, and entry~002 would give the full BSD formula for $E$. The
catalogue itself advertises only the consequence for a density-one set of twists, with entry~006~\cite{OAI006};
the pointwise statement for analytic rank $0$ is our reading of the abstract above.

\begin{proposition}[on paper; conditional on entry~002]\label{prop:sha}
Let $n$ be squarefree and Tunnell's test false. If entry~002 is correct, then
\[
  |\Sha(E_n)|=\Bigl(\frac{A_n-2B_n}{2\,\tau(n)}\Bigr)^2\ \ (n\text{ odd}),\qquad
  |\Sha(E_n)|=\Bigl(\frac{C_n-2D_n}{2\,\tau(n/2)}\Bigr)^2\ \ (n\text{ even}),
\]
where $\tau$ counts divisors. So $|\Sha(E_n)|$ is computed in time $n^{o(1)}$.
\end{proposition}

\begin{proof}
In Tunnell's notation the right sides are $(a(n)/\sigma_0(n))^2$ and $(b(n/2)/\sigma_0(n/2))^2$, which is
Tunnell's conjecture~\cite[p.~330]{Tun83}. Tunnell's remark there computes the local data of $E_n$
($c_p=4$ for odd $p\mid n$, $c_2\in\{2,4\}$, $|E_n(\Q)_{\mathrm{tors}}|=4$) and identifies his conjecture with
the BSD formula, given \cref{cit:tunnell}. By the paragraph above, entry~002 gives that formula. For the cost,
$\tau(m)=\sum_{j\le m}\tau(j)-\sum_{j\le m-1}\tau(j)$ is two runs of the machine \lean{divisorSum_TM_rate} of our
class-number note~\cite{TulCN26}, which is machine-checked; factoring~\cite{Shu26,TulFac26} gives it as well.
\end{proof}

\noindent We checked the formula against the analytic order of $\Sha$ that PARI/GP computes from
$L(E_n,1)$, the real period, the torsion and the Tamagawa numbers, for the $73$ squarefree $n\le300$ with
$L(E_n,1)\ne0$; all agree, and $|\Sha|>1$ occurs $27$ times (for example $|\Sha(E_{17})|=4$,
$|\Sha(E_{113})|=16$). The odd part of Tunnell's conjecture was known before, from Rubin's work on curves with
complex multiplication~\cite{Rub91}\nr; we did not re-check its exact scope.

\begin{remark}[General twists]
For the twists of \cref{cor:twist}, entry~002 would give $|\Sha(E^{-D})|=(L/\Omega)\,|E^{-D}(\Q)_{\mathrm{tors}}|^2/\prod_pc_p$.
The Tamagawa numbers at $p\mid D$ need the primes dividing $D$, so factoring enters, in time $\rate$
by~\cite{Shu26,TulFac26}. We have not worked out the local data in general, so we state no theorem.
\end{remark}

Entry~002 does not change the status of the congruent-number problem itself. In \cref{cor:congruent}(c) the
analytic rank is at least $2$, and entry~002 assumes Selmer corank at most $1$; it says nothing there.

% =====================================================================
\section{What was already within reach}\label{sec:new}
% =====================================================================

In this section we say plainly which forms \cref{thm:main} does not make newly computable. Everything here is
on paper, and ``within reach'' means: in time $n^{o(1)}$ given the engines and known methods.

\paragraph{Binary forms.} For a fixed binary form, $r_A(n)$ is computed in time polynomial in $L$ from the
factorisation of $n$, since the class group of the fixed discriminant is fixed and finite. Factoring takes
$\rate$ with the engines~\cite{RR26,Shu26,TulFac26}.

\paragraph{Even rank.} For $k=4$, Jacobi's formula gives $r_{x^2+y^2+z^2+w^2}(n)=8\sigma(n)$ for odd $n$, so
\cref{thm:main}(ii) computes $\sigma(n)$. Given the factorisation, $\sigma(n)$ is easy; conversely,
$\sigma(n)$ yields the factorisation in randomised polynomial time by Bach, Miller and Shallit~\cite{BMS86}\nr.
So the engine is at least as strong as factoring, which is consistent with~\cite{RR26}. More generally, for even
$k$ the theta series is an Eisenstein series, whose coefficients are twisted divisor sums and need only the
factorisation, plus a cusp form of integral weight. For weight $2$ the cusp coefficients at primes come from point
counts on the associated abelian varieties; for higher weight, Edixhoven, Couveignes and coauthors compute
Ramanujan's $\tau(p)$ in probabilistic polynomial time~\cite{EC11}. So even-rank forms reduce essentially to
factoring and known methods; for weight above $2$ this is modulo the reach of such methods, which we did not check
for the cusp forms that occur here, and modulo their being probabilistic. Our theorem adds a deterministic route there, not new reach.

\paragraph{Forms alone in their genus.} For $x^2+y^2+z^2$, Gauss's formula gives, for squarefree $n\ge5$,
$r_3(n)=12h(-4n)$ for $n\equiv1,2,5,6\pmod8$, $24h(-n)$ for $n\equiv3\pmod8$, and $0$ for $n\equiv7\pmod8$
(we checked it for squarefree $5\le n<800$ against counts of primitive reduced forms). So
\cref{thm:main}(ii) with $a=(1,1,1)$ computes these class numbers in time $n^{o(1)}$. Our class-number
note~\cite{TulCN26} already computes $h(-N)$ in the same time, through reduced forms, so the two results are
consistent and neither adds reach to the other. For a ternary form alone in its genus, Siegel's theorem
expresses $r_A(n)$ through the genus, whose coefficients are, up to local factors that depend on the
factorisation of $n$, class numbers of imaginary quadratic orders, as Gauss's formula illustrates; we state this
only as a pointer, without the exact local factors.

\paragraph{Odd rank: the new part.} For a ternary form with more than one class in its genus, the theta series
also has a cusp part of weight $3/2$. By Waldspurger's theorem\nr, of which \cref{cit:tunnell,cit:bsp} are
explicit cases, the squares of its coefficients at squarefree $n$ are, up to fixed factors, central values of
twisted $L$-functions. We know
of no method that computes these in time $n^{o(1)}$ for a single $n$, even with the engines, other than the
one presented here. This is where \cref{thm:main} adds reach; \cref{sec:tunnell,sec:twists} are examples.

% =====================================================================
\section{The Lean development}\label{sec:formal}
% =====================================================================

In this section we list the modules and state the trusted base. The development is in the public
repository~\cite{TulsyanRepo26}: directories \texttt{lean/EthFollowup/ModP} (pure Mathlib; the files compile
with Lean v4.33.1 and v4.34.1) and \texttt{lean-bridge-anthropic/bridge/EthFollowup/ModP} (on top of \A's
development). \Cref{tab:modules} lists the new modules.

\begin{table}[ht]
\centering\small
\begin{tabular}{@{}p{0.20\textwidth}p{0.66\textwidth}r@{}}
\toprule
module & content & lines\\
\midrule
\lean{RepQF} & \lean{qf}, \lean{repNum}, Cauchy--Schwarz, the adjugate box, the shift, \lean{tForm}, \lean{tunnellTest} & 318\\
\lean{SlpDesc} & straight-line-program descriptions, the canonical run, its count & 337\\
\lean{RepDesc} & monomials, the sign split \eqref{eq:split}, the accumulators, \cref{prop:count} & 600\\
\midrule
\lean{CntPad} & padding of counting packages; the padded rate theorem & 244\\
\lean{RepPoly} & polynomial time of the descriptions & 110\\
\lean{RepPkg} & the family package, \cref{lem:pad}, query specifications, the machines & 320\\
\lean{FinalRep} & the three final theorems & 76\\
\midrule
total & & 2005\\
\bottomrule
\end{tabular}
\caption{The new modules, with line counts (\texttt{wc -l}). The first three are pure Mathlib. The development
reuses, unchanged, the ArithNet, package and machine modules of~\cite{TulCN26} (\texttt{ArithDesc},
\texttt{ArithParams}, \texttt{ArithNet}, \texttt{ArithPackage}, \texttt{CntPackage}, \texttt{GRate},
\dots), the factoring modules of~\cite{TulFac26} and the bridge modules of~\cite{Tul26}.}\label{tab:modules}
\end{table}

\paragraph{Trusted base and axioms.} \texttt{\#print axioms} reports \lean{propext}, \lean{Classical.choice}
and \lean{Quot.sound} for the three theorems of \cref{thm:main}; we ran it on 5~November 2026 against a build of
the pinned development whose sources are identical to the files of \cref{tab:modules}. No new module contains
\lean{sorry}, a new axiom or \lean{native_decide}.
%PENDING-REPLAY
A fresh \lean{leanchecker} replay of the dependency closure of \lean{EthFollowup.ModP.FinalRep}, the module
holding the three final theorems, is pending at the time of writing.
What must be trusted is the Lean kernel, the Mathlib definitions named in \cref{sec:statement}, our definitions
quoted there, and \A's definition of \texttt{Outputs}. \A's development, Mathlib's proofs and our other modules
are checked by the kernel; they are part of the checked closure, not of the trusted base.

\paragraph{One design point.} The padding is generic: \lean{CntPackage.padded} works for any counting
package whose rounds admit the original query words, and the rate theorem for it is proved once. Padding the
input word, instead of enlarging the box inside the description family, leaves ArithNet and its parameter proofs
untouched.

% =====================================================================
\section{Related work}\label{sec:related}
% =====================================================================

We organise related work by idea: computing representation numbers and Tunnell's counts, central values of
twists, the congruent-number problem, and the counting engines with their formal proofs.

\paragraph{Representation numbers and Tunnell's counts.} Tunnell noted the $O(n^{3/2})$ cost of his
criterion~\cite[p.~325]{Tun83}. Counting representations of $n-cz^2$ by binary forms brings one $n$ to about
$n^{1/2+o(1)}$, and Hart, Tornar\'ia and Watkins put the cost of this route up to $N$ at
$O(N^{3/2})$~\cite[\S1.3]{HTW10}. They computed Tunnell's counts for all $n\le10^{12}$ by out-of-core FFT
multiplication of theta series, split by residue classes following Elkies, in about $30$ hours for the classes
$1,3\pmod8$ on a $16$-core machine~\cite{HTW10}. That is quasi-linear in the range, so $X^{o(1)}$ per coefficient
on average, but with memory about $X$ and no gain for one $n$. In contrast, we compute one $n$ in $n^{o(1)}$, for
every fixed positive-definite form, at the price of galactic constants.

\paragraph{Central values of twists.} The approximate functional equation computes $L(E\otimes\chi_D,1)$
numerically in time ``very roughly proportional to $|D|$''~\cite{MRVT07}. Mao, Rodriguez-Villegas and
Tornar\'ia compute the coefficients of the theta combination for all $|D|\le x$ in about $x^{3/2}$ steps, by
running over the lattice points of ellipsoids, and extend the method to real twists and to $L(E,1)=0$ with
generalised theta series~\cite{MRVT07}; Pacetti and Tornar\'ia treat composite levels~\cite{PT08}. Closest to
our work are these papers, whose formulas we use. In contrast, we evaluate one coefficient in $|D|^{o(1)}$ instead
of enumerating the ellipsoid.

\paragraph{The congruent-number problem.} Tunnell's criterion~\cite{Tun83} is unconditional in one direction by
Coates and Wiles~\cite{CW77}\nr; Conrad gives an exposition~\cite{Con08}. For $n\equiv5,6,7\pmod8$ much is
known unconditionally; for example, Smith proves that at least $55.9\%$ of squarefree $n\equiv5,6,7\pmod8$ are
congruent~\cite{Smi16}. OpenAI's entries~002 and~006~\cite{OAI002,OAI006} give, if correct, the full BSD formula in
Selmer corank at most $1$ and Goldfeld's density statements for twists; neither decides the analytic-rank-$\ge2$
case of \cref{cor:congruent}(c). In contrast, our contribution to this problem is the speed of the test, not its
logical status.

\paragraph{Counting engines and formal proofs.} \Oa~\cite{OpenAI26} and \A~\cite{Anthropic26} claim deterministic
subexponential $\#$SAT algorithms, with Lean developments, under review. Rudich and Rousseau~\cite{RR26} factored
in $N^{o(1)}$ with them, Shunia~\cite{Shu26} proved the rate $\rate$ on paper, and~\cite{IF26} and our factoring
note~\cite{TulFac26} proved it in Lean. Our earlier notes applied \A's engine to Hamiltonian cycles~\cite{Tul26},
Eulerian circuits~\cite{TulEuler26}, class numbers and lattice-point counts~\cite{TulCN26}, and the parity of
$\pi(x)$~\cite{TulPar26}. The closest is the class-number note, whose library we reuse; in contrast, the
equations here have $\Theta(k^2)$ products, integer unknowns and values of $2L$ bits, which the shift, the sign
split and the padding handle. We use Lean~4~\cite{Lean4} and Mathlib~\cite{Mathlib20}; Mathlib at the commit we use
proves Lagrange's four-square theorem, and we found in it no representation counts and no congruent numbers.

\paragraph{Novelty search.} On 5~November 2026 we grepped the titles and abstracts of all $716$ works on Hexagon
through its public query interface for representation numbers, ternary and quadratic forms, theta series,
congruent numbers, Tunnell, Waldspurger, BSD, central values and elliptic curves; queried arXiv's interface for the
Exponential Time Hypothesis or $\#$SAT together with quadratic forms, congruent numbers and Tunnell; and grepped
OpenAI's catalogue for the same terms. This complements our novelty audit of the same day~\cite{TulsyanRepo26}
(file \texttt{research/novelty-audit-2026-11-05.md}). We found no work that computes representation numbers,
Tunnell's counts or these $L$-values with either ETH engine. A negative search result is weak evidence.

% =====================================================================
\section{Limitations}\label{sec:limitations}
% =====================================================================

\begin{itemize}
\item \emph{The results stand or fall with \A's development.} Every theorem rests on Anthropic's Lean development
of its claimed refutation of the Exponential Time Hypothesis, and the labs' claims are under review. Lean checks
that development as part of our proof, and our statements use only Mathlib, our quoted definitions and \A's
\texttt{Outputs}. We do need the development, built from source at the pinned commit, to be sound, and the
toolchain to be sound. The fresh kernel replay of \lean{EthFollowup.ModP.FinalRep} is pending.
\item \emph{The consequences are on paper only.} Lean proves that the machines output counts. The link to
congruent numbers, $L$-values and Tate--Shafarevich groups (\cref{sec:tunnell,sec:twists,sec:sha}) rests on cited
theorems, some of which we did not re-read (marked~\nr), and is not formalised. \Cref{cor:congruent}(c) needs BSD.
\item \emph{\Cref{sec:sha} is conditional on an unrefereed result.} OpenAI's entry~002 is a set of preprints that
we have not checked.
\item \emph{The twist corollary has limited scope.} It needs $L(E,1)\ne0$, squarefree conductor and the residue
conditions of \cref{cor:twist}, and the constant $\kappa_E$ is fixed but not computed by our machines. Real twists
need weighted counts that \cref{thm:main} does not cover.
\item \emph{The algorithms are galactic.} \A's levels start at $K_0=2^{x_0}$ with $x_0\ge18$, the level needs
$K^{40}\lesssim s$, and the constants are astronomical: for $k=3$ the shape constant alone is
$64\cdot299\cdot148=2{,}832{,}128$, before the factor $2^{12}$ of the package. The machines are of no practical use and
beat $n^{1/2}$ only for astronomically large $n$.
\item \emph{The machines use exponential space.}
\item \emph{No human has checked the work.} The Lean code and this text were written with AI assistance. No human
expert has yet checked that \cref{sec:statement} lists everything a reader must trust, or the derivations on paper.
\end{itemize}

\section{Conclusion}\label{sec:conclusion}
We have shown that a representation number of a fixed positive-definite form is one count of a straight-line
program with nonnegative coefficients, after a shift and a sign split, and that the counting engine evaluates it in
$n^{o(1)}$ steps once the input is padded. The arithmetic meaning is classical: Tunnell's counts give $L(E_n,1)$, and
ternary theta combinations give central values of twists. The result is a step towards, not a practical method for,
these computations. An interesting direction is whether the same route handles the weighted counts of
generalised theta series, which would cover real twists and base curves of positive rank, and whether a count
can certify $L'(E_n,1)\ne0$, which would make the test informative for $n\equiv5,6,7\pmod8$.

\section*{Acknowledgements}
This work was done with AI assistance (Claude Code, Anthropic): the constructions, the Lean development, the
numerical checks, the literature and citation checks and the text were produced in sessions with Claude Code. The
Lean theorems are checked by the Lean kernel; everything else awaits human expert review. Responsibility for any
errors rests with the author. This note continues the project of~\cite{Tul26,TulFac26,TulCN26,TulPar26}.

\begin{thebibliography}{99}

\bibitem{Anthropic26}
Claude (Anthropic), \emph{A $2^{o(n)}$ time algorithm for {3-SAT}}, Hexagon preprint hexagon:2610.00577v1,
October 2026, \url{https://hexagonmath.org/abs/2610.00577v1}. Lean formalisation in the repository
\url{https://github.com/anthropics/formal-math}, folder \texttt{eth}; we use commit
\texttt{16a5ce41304fcdaea4e0be7c132e1c361dc26da5}.

\bibitem{BMS86}
E.~Bach, G.~Miller, and J.~Shallit, \emph{Sums of divisors, perfect numbers and factoring}, SIAM J. Comput.
\textbf{15} (1986), 1143--1154.

\bibitem{BSP90}
S.~B\"ocherer and R.~Schulze-Pillot, \emph{On a theorem of {W}aldspurger and on {E}isenstein series of {K}lingen
type}, Math. Ann. \textbf{288} (1990), 361--388.

\bibitem{CW77}
J.~Coates and A.~Wiles, \emph{On the conjecture of {B}irch and {S}winnerton-{D}yer}, Invent. Math. \textbf{39}
(1977), 223--251.

\bibitem{Con08}
K.~Conrad, \emph{The congruent number problem}, Harvard College Math. Review \textbf{2} (2008).

\bibitem{EC11}
B.~Edixhoven and J.-M.~Couveignes (eds.), \emph{Computational aspects of modular forms and {G}alois
representations}, Annals of Math. Studies, vol.~176, Princeton Univ. Press, 2011; arXiv:math/0605244.

\bibitem{Gro87}
B.~H.~Gross, \emph{Heights and the special values of {$L$}-series}, Number Theory (Montreal, 1985), CMS Conf.
Proc., vol.~7, Amer. Math. Soc., Providence, RI, 1987, pp.~115--187.

\bibitem{GZ86}
B.~H.~Gross and D.~B.~Zagier, \emph{Heegner points and derivatives of {$L$}-series}, Invent. Math. \textbf{84}
(1986), 225--320.

\bibitem{HTW10}
W.~Hart, G.~Tornar\'ia, and M.~Watkins, \emph{Congruent number theta coefficients to $10^{12}$}, Algorithmic
Number Theory (ANTS~IX), Lecture Notes in Comput. Sci., vol.~6197, Springer, 2010, pp.~186--200.

\bibitem{IF26}
1337isnot1337 (GitHub user) and contributors, \emph{integer-factoring: deterministic integer factoring bounds},
public repository \url{https://github.com/1337isnot1337/integer-factoring}; Lean formalisation in
\texttt{formal-appendix/} (pull request \#20, merged 30~October 2026).

\bibitem{Kol88}
V.~A.~Kolyvagin, \emph{Finiteness of {$E(\mathbb Q)$} and {$\Sha(E,\mathbb Q)$} for a subclass of {W}eil curves},
Izv. Akad. Nauk SSSR Ser. Mat. \textbf{52} (1988), no.~3, 522--540.

\bibitem{Kol90}
V.~A.~Kolyvagin, \emph{Euler systems}, The Grothendieck Festschrift, Vol.~II, Birkh\"auser, 1990.

\bibitem{Lean4}
L.~de Moura and S.~Ullrich, \emph{The {L}ean 4 theorem prover and programming language}, Automated Deduction --
CADE 28, Lecture Notes in Comput. Sci., vol. 12699, Springer, 2021, pp.~625--635.

\bibitem{MRVT07}
Z.~Mao, F.~Rodriguez-Villegas, and G.~Tornar\'ia, \emph{Computation of central value of quadratic twists of
modular {$L$}-functions}, Ranks of Elliptic Curves and Random Matrix Theory, London Math. Soc. Lecture Note
Ser., vol.~341, Cambridge Univ. Press, 2007, pp.~273--288; arXiv:math/0606762.

\bibitem{Mathlib20}
The mathlib Community, \emph{The {L}ean mathematical library}, Proc. 9th ACM SIGPLAN International Conference on
Certified Programs and Proofs (CPP 2020), 2020, pp.~367--381.

\bibitem{OAI002}
OpenAI, \emph{Exact {B}irch--{S}winnerton-{D}yer formula from low {S}elmer corank}, preprint dated 7~October 2026,
with companion preprints of 6--7~October 2026; entry~002, ``The full BSD formula from low Selmer corank'', of the
catalogue \texttt{CONTENTS.md} in the repository \url{https://github.com/openai/math}. Not refereed.

\bibitem{OAI006}
OpenAI, entry~006, ``Goldfeld's conjecture: densities and mean analytic rank'', of the catalogue
\texttt{CONTENTS.md} in the repository \url{https://github.com/openai/math}, 2026. Not refereed.

\bibitem{OpenAI26}
OpenAI, \emph{The {E}xponential {T}ime {H}ypothesis is false}, preprint dated 17~October 2026, revised 25~October
2026. Paper and Lean formalisation in the repository \url{https://github.com/openai/math}.

\bibitem{PT08}
A.~Pacetti and G.~Tornar\'ia, \emph{Computing central values of twisted {$L$}-series: the case of composite
levels}, Experiment. Math. \textbf{17} (2008), 459--472; arXiv:math/0607008.

\bibitem{Pal12}
V.~Pal, \emph{Periods of quadratic twists of elliptic curves}, with an appendix by A.~Agashe, Proc. Amer. Math.
Soc. \textbf{140} (2012), 1513--1525; arXiv:1012.0094.

\bibitem{PARI}
The PARI~Group, \emph{PARI/GP version 2.17.2}, Univ. Bordeaux, \url{https://pari.math.u-bordeaux.fr/}.

\bibitem{RR26}
I.~Rudich and L.-M.~Rousseau, \emph{Deterministic factoring in time $N^{o(1)}$}, Hexagon preprint
hexagon:2610.00600v1, October 2026, \url{https://hexagonmath.org/abs/2610.00600v1}.

\bibitem{Rub91}
K.~Rubin, \emph{The ``main conjectures'' of {I}wasawa theory for imaginary quadratic fields}, Invent. Math.
\textbf{103} (1991), 25--68.

\bibitem{Shu26}
J.~Shunia, \emph{An improved subexponential bound for deterministic integer factorization}, Hexagon preprint
hexagon:2610.00711v1, 2026, \url{https://hexagonmath.org/abs/2610.00711v1}.

\bibitem{Smi16}
A.~Smith, \emph{The congruent numbers have positive natural density}, arXiv:1603.08479, 2016.

\bibitem{Tun83}
J.~B.~Tunnell, \emph{A classical {D}iophantine problem and modular forms of weight {$3/2$}}, Invent. Math.
\textbf{72} (1983), 323--334.

\bibitem{Tul26}
V.~Tulsyan, \emph{Counting {H}amiltonian cycles in sparse graphs in subexponential time, assuming the new
$\#${SAT} algorithms}, Hexagon preprint hexagon:2610.00672v1, October 2026,
\url{https://hexagonmath.org/abs/2610.00672v1}.

\bibitem{TulCN26}
V.~Tulsyan, \emph{Class numbers and lattice-point counts in time $2^{O(n/\sqrt{\log n})}$, assuming the new
$\#${SAT} algorithms}, Hexagon preprint hexagon:2611.00115v2, November 2026,
\url{https://hexagonmath.org/abs/2611.00115}.

\bibitem{TulEuler26}
V.~Tulsyan, \emph{Counting {E}ulerian circuits in subexponential time, assuming the new $\#${SAT} algorithms},
Hexagon preprint hexagon:2611.00088v1, November 2026, \url{https://hexagonmath.org/abs/2611.00088v1}.

\bibitem{TulFac26}
V.~Tulsyan, \emph{An independent machine-checked proof of deterministic factoring in time
$2^{O(n/\sqrt{\log n})}$, via the new $\#${SAT} algorithms}, preprint, November 2026; source in the directory
\texttt{factoring-note/} of~\cite{TulsyanRepo26}.

\bibitem{TulPar26}
V.~Tulsyan, \emph{The parity of the prime-counting function in time $x^{o(1)}$, assuming the new $\#${SAT}
algorithms}, Hexagon preprint hexagon:2611.00116v1, November 2026, \url{https://hexagonmath.org/abs/2611.00116v1}.

\bibitem{TulsyanRepo26}
V.~Tulsyan, \emph{eth-hamiltonicity: {L}ean 4 formalisation and research notes}, public repository
\url{https://github.com/xennygrimmato/eth-hamiltonicity}, 2026; Lean sources as of commit \texttt{3e181e7}.

\end{thebibliography}

% =====================================================================
% CITATION CHECK (5 November 2026).
%  Tun83   Scan read (GDZ PPN356556735_0072): Main Theorem p. 323, cost remark p. 325, Theorem 3 pp. 328-329,
%          root-number remark p. 329, Theorem 4, Conjecture and Remark p. 330. Theorem 3 re-checked numerically
%          against PARI for 183 squarefree n <= 300.
%  BSP90   Scan read (GDZ PPN235181684_0288): setup sec. 1.4 p. 365, Theorem (3.5) p. 378, remarks p. 378.
%  MRVT07  arXiv math/0606762 read: (3.1) attributing l = 1 to Gross Prop. 13.5; forms (4.1), (4.2); cost remarks.
%  PT08    arXiv math/0607008 read: quotation on the fraction of twists covered.
%  Pal12   arXiv:1012.0094 read: Main Result 1.1, Prop. 2.5, Cor. 2.6; remark on misquotation.
%  HTW10   WRAP preprint read: sec. 1.3 O(N^{3/2}); 10^12 timings.
%  Con08   Theorem 15 read.
%  Gro87   Not obtained; used through MRVT07 (3.1).
%  OAI002, OAI006: local copy of openai/math CONTENTS.md; manuscript abstract quoted from the entry.
%  Smi16   arXiv abstract (55.9% density statement).
%  EC11    arXiv math/0605244 abstract; probabilistic polynomial time.
%  Marked \dag (not re-read for this note): CW77 (quoted by Tunnell p. 329), GZ86, Kol88, Kol90, Rub91, BMS86;
%          rationality of L/Omega for twists.
%  Anthropic26: Hexagon API lists hexagon:2610.00577v1 (author Anthropic, 27 Oct 2026). Shu26: API lists
%          hexagon:2610.00711v1. Hexagon IDs of our notes from barrier-note/main.tex and the API.
%  Novelty search (5 Nov 2026, ~20:35 UTC): Hexagon API, 716 entries, regex over titles and abstracts; arXiv API
%          queries returned no hits; OpenAI CONTENTS.md local copy grepped: no hits.
% =====================================================================

\end{document}
